\title{DeepSeekMath: Pushing the Limits of Mathematical Reasoning in Open Language Models}

\begin{abstract}
The task of mathematical reasoning remains particularly demanding for language models owing to its inherent complexity and structured requirements. In this study, we present DeepSeekMath~7B, an extension of DeepSeek-Coder-Base-v1.5 7B further pre-trained with 120B math-specific tokens curated from Common Crawl alongside associated natural language and code resources. Without recourse to external toolkits or ensembling approaches, DeepSeekMath~7B attains a notable 51.7\% score on the rigorous MATH benchmark, closing the gap with models such as Gemini-Ultra and GPT-4. When evaluated using self-consistency across 64 samples, DeepSeekMath~7B achieves a score of 60.9\% on MATH. Two principal strategies underlie the mathematical reasoning advances in DeepSeekMath~7B: First, we leverage the untapped promise of publicly available internet data through a carefully crafted pipeline for data selection. Second, we introduce Group Relative Policy Optimization (GRPO), a specialized variant of Proximal Policy Optimization (PPO) that simultaneously improves mathematical problem-solving and enhances the memory efficiency of PPO.
\end{abstract}

\section{Introduction}

Large language models (LLM) have transformed the landscape of mathematical reasoning in artificial intelligence, driving substantial improvements across benchmarks for quantitative reasoning~\citep{MATH} as well as geometry reasoning~\citep{trinh2024solving}. In addition to these advances, such models increasingly serve as valuable tools in aiding humans to tackle intricate mathematical challenges~\citep{tao}. Nevertheless, state-of-the-art systems like GPT-4~\citep{gpt4} and Gemini-Ultra~\citep{gemini} remain closed, while the performance of current open-source models still lags significantly behind.

In this paper, we present DeepSeekMath, a specialized language model that exhibits superior mathematical proficiency compared to open-source alternatives and closely matches GPT-4’s results on standard academic evaluations. The foundation of our approach is the DeepSeekMath Corpus, a comprehensive and high-quality pre-training dataset consisting of 120B math-specific tokens. This corpus is assembled from Common Crawl (CC) data, identified using a fastText classifier \citep{joulin2016fasttext}. Initially, the classifier is trained with positive samples from OpenWebMath \citep{openwebmath} alongside a heterogeneous mix of web pages as negative samples. We then leverage the classifier to extract more math-related content from CC, which is subsequently refined through human validation. The enriched set is used to retrain the classifier, further optimizing its precision. According to our evaluations, the resulting dataset is of exceptional quality: the DeepSeekMath-Base 7B model attains 64.2\% accuracy on GSM8K \citep{gsm8k} and 36.2\% on the advanced MATH dataset \citep{MATH}, surpassing the performance of Minerva 540B \citep{minerva}. Furthermore, because the DeepSeekMath Corpus encompasses multiple languages, we observe notable gains on Chinese mathematical assessments \citep{wei2023cmath,agieval}. Our work in curating and processing mathematical data offers valuable insights for the research community and highlights considerable opportunities for future advancements.

DeepSeekMath-Base is built upon the DeepSeek-Coder-Base-v1.5 7B \citep{deepseek-coder}, as our findings suggest that initializing from a code-oriented pretrained model offers notable advantages over general-purpose LLMs. Additionally, we find that mathematical pretraining not only strengthens the model’s mathematical competency, but also leads to improvements on broader reasoning benchmarks, including MMLU \citep{mmlu} and BBH \citep{bbh}.

Following pretraining, we subject DeepSeekMath-Base to mathematical instruction tuning, incorporating datasets featuring chain-of-thought \citep{cot}, program-of-thought \citep{pot,pal}, and tool-augmented reasoning \citep{tora}. The final model, DeepSeekMath-Instruct 7B, outperforms all other models of equivalent size and even rivals instruction-tuned models of 70B parameters from the open-source community.

We also propose Group Relative Policy Optimization (GRPO), a reinforcement learning approach derived from Proximal Policy Optimization (PPO) \citep{schulman2017proximal}. Unlike PPO, GRPO eliminates the need for a separate critic network and instead computes baselines using group-level scores, thereby making training more computationally efficient. Using only a fraction of the English instruction tuning corpus, GRPO achieves significant performance gains over DeepSeekMath-Instruct. These gains are evident in both in-domain tasks (GSM8K: 82.9\% $\rightarrow$ 88.2\%, MATH: 46.8\% $\rightarrow$ 51.7\%) and out-of-domain evaluations (e.g., CMATH: 84.6\% $\rightarrow$ 88.8\%) during reinforcement learning. Moreover, we articulate a unified framework that connects various methods, such as Rejection Sampling Fine-Tuning (RFT) \citep{yuan2023scaling}, Direct Preference Optimization (DPO) \citep{dpo}, PPO, and our GRPO, demonstrating that these techniques can be interpreted as manifestations of either direct or simplified reinforcement learning paradigms. Through extensive empirical studies—comparing, for instance, online versus offline training, outcome-based versus process-oriented supervision, and single-turn versus iterative reinforcement learning—we probe the fundamental properties underlying this unified framework. Finally, we offer explanations for the observed performance improvements of instruction-tuned models under RL, and outline promising research directions for advancing reinforcement learning using insights gained from this unified perspective.

\subsection{Contributions}

This work advances both large-scale math-specific pre-training and the systematic study of reinforcement learning techniques.

\textbf{Math Pre-Training at Scale}

\begin{itemize}[topsep=0pt]
    \item We demonstrate that Common Crawl, a publicly available internet dataset, holds significant mathematical content. Through a carefully crafted data selection process, we assemble the DeepSeekMath Corpus—a high-quality, 120B-token dataset extracted from web resources selected for mathematical relevance. This resource substantially exceeds the scale of previous efforts, being almost sevenfold larger than Minerva’s mathematical subset \citep{minerva} and about nine times greater than OpenWebMath’s initial release \citep{openwebmath}.

    \item Our DeepSeekMath-Base 7B model, after pre-training, rivals the performance of Minerva 540B \citep{minerva}, suggesting that model size is not the sole determinant of proficiency in mathematical reasoning. High-quality, domain-specific pre-training allows comparatively smaller models to attain strong results.

    \item We present insights derived from training experiments on mathematics. Sequential training—first on code, then on math tasks—enhances the model’s mathematical problem-solving abilities, regardless of tool assistance. This result provides partial empirical support for the hypothesis: \textit{does exposure to code enhance reasoning skills?} Our evidence affirms this for the domain of mathematical reasoning.

    \item Despite the prevalent practice of training on arXiv content in mathematics-focused research, such pre-training yields no observable benefit across the mathematical evaluation benchmarks examined in this study.
\end{itemize}

\textbf{Exploration and Analysis of Reinforcement Learning}

\begin{itemize}[topsep=0pt]
    \item We propose Group Relative Policy Optimization (GRPO), a reinforcement learning approach characterized by both its efficiency and effectiveness. GRPO eliminates the need for a critic by utilizing group-based score estimations as the baseline, thereby markedly decreasing training costs in comparison to Proximal Policy Optimization (PPO).
    \item Our empirical results reveal that applying GRPO yields substantial improvements in the instruction-tuned model DeepSeekMath-Instruct, relying exclusively on instruction-tuning data. Additionally, we note improvements in out-of-domain performance during the course of reinforcement learning.
    \item We establish a comprehensive framework to contextualize and compare methodologies such as RFT, DPO, PPO, and GRPO. We also undertake an extensive suite of experiments—exploring, for example, online versus offline learning, outcome versus process-based supervision, and single-turn compared to iterative reinforcement learning—to rigorously analyze the foundational components underpinning this framework.
    \item Utilizing this integrated paradigm, we delve into the underlying mechanisms responsible for the effectiveness of reinforcement learning, and outline several promising directions for advancing reinforcement learning techniques for LLMs.
\end{itemize}

\subsection{Summary of Evaluations and Metrics}

\begin{itemize}[topsep=0pt]
    \item \textbf{English and Chinese Mathematical Reasoning}: Our models are systematically evaluated using a diverse array of English and Chinese datasets that encompass mathematical problems from elementary to undergraduate levels. English evaluation sets feature GSM8K \citep{gsm8k}, MATH \citep{MATH}, SAT \citep{llemma}, OCW Courses \citep{minerva}, and MMLU-STEM \citep{mmlu}. For Chinese, we utilize MGSM-zh \citep{mgsm}, CMATH \citep{wei2023cmath}, Gaokao-MathCloze \citep{agieval}, and Gaokao-MathQA \citep{agieval}. The assessments consider both the capability of models to produce self-contained, tool-free solutions, as well as their proficiency in leveraging Python for problem-solving.

    Across the English benchmarks, DeepSeekMath-Base attains competitive results relative to the proprietary Minerva 540B \citep{minerva}, outperforming all existing open-source base models—including Mistral 7B \citep{mistral} and Llemma-34B \citep{llemma}—irrespective of their exposure to math-specific pre-training, and often by a notable extent. On the Chinese datasets, DeepSeekMath-Base exhibits further superiority, attributed in part to our decision to extend pre-training to encompass high-quality mathematical data in languages beyond English, diverging from prior approaches \citep{minerva,llemma}. Augmenting with instruction tuning and reinforcement learning, both DeepSeekMath-Instruct and DeepSeekMath-RL display robust results—most prominently, they achieve greater than 50\% accuracy on the MATH competition benchmark, representing a first among open-source initiatives.
\end{itemize}

\item \textbf{Formal Mathematics}: We assess DeepSeekMath-Base on the informal-to-formal theorem proving task from \citep{dsp_proof}, employing the miniF2F benchmark \citep{minif2f} with Isabelle \citep{isabelle} as the designated proof assistant. The results indicate that DeepSeekMath-Base achieves notable few-shot autoformalization performance.
\item \textbf{Natural Language Understanding, Reasoning, and Code}: To provide a holistic evaluation of the model’s general language understanding, reasoning, and programming abilities, we measure DeepSeekMath-Base’s performance on the Massive Multitask Language Understanding (MMLU) benchmark \citep{mmlu}, which spans 57 diverse multiple-choice tasks, as well as BIG-Bench Hard (BBH) \citep{bbh}, a collection of 23 complex tasks primarily involving multi-step reasoning. Additionally, we include assessments on HumanEval \citep{codex} and MBPP \citep{mbpp}, standard benchmarks for code generation models. The results show that math-focused pre-training improves both reasoning and language comprehension metrics.

\section{Math Pre-Training}

\subsection{Data Collection and Decontamination} Here, we describe the methodology used to assemble the DeepSeekMath Corpus from Common Crawl. As illustrated in Figure \ref{fig:data_collect}, we devised an iterative pipeline that systematically acquires a large-scale mathematical dataset, beginning with a seed collection, such as a small, high-quality mathematical corpus. Notably, this framework can be generalized to other specialized domains, for example, coding.

Our initial seed corpus is OpenWebMath \citep{openwebmath}, consisting of curated, high-quality mathematical web resources. Using this dataset, we train a fastText model \citep{joulin2016fasttext} to retrieve additional web pages similar in content and quality to OpenWebMath. For training, we sample 500,000 entries from the seed corpus as positive examples, alongside 500,000 negative samples randomly drawn from Common Crawl. We utilize an open-source implementation\footnote{\url{https://fasttext.cc}} of fastText, with hyperparameters set as follows: vector dimensionality of 256, learning rate of 0.1, maximum word n-gram length of 3, minimum word frequency of 3, and training for 3 epochs. To manage the scale of Common Crawl, we first apply both URL-based deduplication and near-duplicate filtering, yielding a collection of 40B unique HTML pages. The fastText model is then employed to select mathematical pages from this deduplicated pool. To exclude lower-quality materials, all retrieved web pages are ranked based on their fastText-assigned scores, retaining only those with the highest rankings. We determine how much data to retain by conducting pre-training experiments across top 40B, 80B, 120B, and 160B tokens. In the initial iteration, only the top 40B tokens are selected for further use.

Following the initial round of data acquisition, a significant portion of mathematical web pages remained uncaptured. This shortfall primarily stems from the limited diversity in the positive samples used to train the fastText model. To address this, we seek out further mathematical websites to broaden and enhance the seed set, aiming to refine the fastText model’s performance. Concretely, we first partition the entirety of Common Crawl into non-overlapping domains, where each domain is defined as the collection of web pages sharing a common base URL. For every domain, we determine the fraction of its web pages included in the initial data collection pass. If over 10\% of a domain’s pages have already been obtained, we categorize it as being related to mathematics (for instance, \url{mathoverflow.net}). Next, we undertake manual annotation to pinpoint specific URLs within these domains that pertain to mathematical content (e.g., \url{mathoverflow.net/questions}). Pages corresponding to these URLs that were not previously collected are incorporated into the expanded seed set. Through this method, we increase the diversity of positive samples and thereby retrain the fastText model to capture a broader array of mathematical data in subsequent iterations. Ultimately, after four rounds of collection, we assemble a dataset comprising 35.5M math-oriented web pages, containing a total of 120B tokens. We observe that, by the fourth pass, nearly 98\% of the pages obtained had already been gathered by the third round, leading us to conclude the data collection process at this point.

To guard against benchmark contamination, we adopt the filtering strategy from \cite{deepseek-coder}, systematically excluding any web content that replicates questions or answers from widely used mathematical benchmarks. This encompasses English datasets such as GSM8K \citep{gsm8k} and MATH \citep{MATH}, as well as Chinese benchmarks like CMATH \citep{wei2023cmath} and AGIEval \citep{agieval}. The filtering rule is as follows: any text segment that contains a substring of exactly 10 tokens, which matches any portion from these benchmarks, is omitted from the training data. In the case of benchmark items shorter than 10 grams but at least 3 grams long, we perform exact string matching to excise potentially contaminated pages from our math corpus.

\subsection{Validating the Quality of the DeepSeekMath~Corpus}

To assess the quality of the DeepSeekMath~Corpus, we conduct a set of pre-training experiments in comparison with several recently introduced math-focused training datasets:

\begin{itemize}[topsep=0pt]
    \item \textbf{MathPile} \citep{mathpile}: This multi-source collection encompasses 8.9B tokens, originating from textbooks, Wikipedia, ProofWiki, CommonCrawl, StackExchange, and arXiv, with more than 85\% sourced from arXiv;
    \item \textbf{OpenWebMath} \citep{openwebmath}: Consists of CommonCrawl web pages filtered for mathematical content, with a total size of 13.6B tokens;
    \item \textbf{Proof-Pile-2} \citep{llemma}: A composite mathematics dataset, including OpenWebMath, AlgebraicStack (10.3B tokens of mathematical code), and arXiv publications (28.0B tokens). Experiments on Proof-Pile-2 follow \cite{llemma}, using a source ratio of 2:4:1 for arXiv, web, and code data, respectively.
\end{itemize}

\subsubsection{Training Setting}
\label{sec:quality-policy}
We conduct specialized math training on a foundational language model containing 1.3B parameters, architecturally consistent with the DeepSeek LLMs \citep{deepseek-llm}, which we refer to as DeepSeek-LLM 1.3B. For each distinct mathematical dataset, we train a separate instance of the model over 150B tokens. All training runs employ the lightweight and efficient HAI-LLM \citep{haillm} framework. Mirroring the approach used in DeepSeek LLMs, we leverage the AdamW optimizer \citep{adamW}, configuring $\beta_1=0.9$, $\beta_2=0.95$, and $\mathrm{weight\_decay}=0.1$. The learning rate follows a multi-step schedule: it ramps up to its peak within 2,000 warmup steps, is reduced to 31.6\% of its maximum by the 80\% completion mark of training, and further lowered to 10.0\% by the 90\% point. The highest learning rate employed is 5.3e-4. Training utilizes a batch size corresponding to 4M tokens, with a context window capped at 4K tokens.

\subsubsection{Evaluation Results}

\textbf{DeepSeekMath~Corpus exhibits exceptional quality, robust multilingual coverage in mathematical content, and surpasses others in size.}

\begin{itemize}[topsep=0pt]
    \item \textbf{High-quality}: Downstream task performance is assessed on eight mathematical evaluation sets with few-shot chain-of-thought prompting \cite{cot}. According to Table \ref{tab:corpora_comparison}, the DeepSeekMath~Corpus-trained model consistently achieves superior results. Figure \ref{fig:corpora_comparisons} further demonstrates that models trained on DeepSeekMath outperform those trained on Proof-Pile-2 at 50B tokens (corresponding to a single epoch over Proof-Pile-2), suggesting that the DeepSeekMath~Corpus maintains higher average data quality.
    \item \textbf{Multilingual}: The DeepSeekMath~Corpus includes mathematical material in various languages, with English and Chinese being the most prominent. Results in Table \ref{tab:corpora_comparison} indicate that training with DeepSeekMath~Corpus yields improved mathematical reasoning in both English and Chinese. By contrast, previous corpora—which largely concentrate on English—display limited benefit or even negatively impact the model’s capability for mathematical reasoning in Chinese.
    \item \textbf{Large-scale}: The volume of DeepSeekMath~Corpus is substantially greater than that of prior mathematical corpora. As visualized in Figure \ref{fig:corpora_comparisons}, training DeepSeek-LLM 1.3B on this corpus yields a steeper learning progression and more persistent gains. Meanwhile, the comparative corpora, due to their smaller size and repeated exposures during training, lead to models whose performance plateaus rapidly.
\end{itemize}

\subsection{Training and Evaluating DeepSeekMath-Base 7B}

This section presents DeepSeekMath-Base 7B, a foundational model designed for advanced mathematical reasoning. The model is built upon DeepSeek-Coder-Base-v1.5 7B \citep{deepseek-coder} and trained on a dataset comprising 500B tokens. The dataset is distributed as follows: 56\% from the DeepSeekMath Corpus, 4\% from AlgebraicStack, 10\% from arXiv, 20\% from Github code repositories, and the remaining 10\% from Common Crawl natural language data in both English and Chinese. The training methodology largely adheres to the protocol outlined in Section \ref{sec:quality-policy}, with the exception that the peak learning rate is increased to 4.2e-4, and a batch size of 10M tokens is employed.

We thoroughly evaluate the mathematical proficiency of DeepSeekMath-Base 7B. This evaluation considers the model’s capabilities in generating comprehensive mathematical solutions independently of external assistance, its performance on mathematical tasks that involve external tools, as well as its capacity for formal theorem proving. In addition to these mathematics-specific tasks, the model’s competencies in natural language comprehension, logical reasoning, and programming are also examined to offer a broader perspective on its general capabilities.

\paragraph{Mathematical Problem Solving with Step-by-Step Reasoning}

To assess DeepSeekMath-Base’s abilities in mathematical problem-solving, we employ few-shot chain-of-thought prompting \citep{cot} across eight evaluation benchmarks in both English and Chinese. These benchmarks span a range of quantitative reasoning tasks (including GSM8K \citep{gsm8k}, MATH \citep{MATH}, and CMATH \citep{wei2023cmath}) as well as multiple-choice assessments (such as MMLU-STEM \citep{mmlu} and Gaokao-MathQA \citep{agieval}), thus representing a wide spectrum of mathematical topics from basic to advanced (elementary through college-level) difficulty.

As indicated in Table \ref{tab:base_math_cot}, DeepSeekMath-Base 7B consistently achieves top results across all eight evaluated benchmarks among open-source base models, which includes both the popular general-purpose Mistral 7B \citep{mistral} and the recently introduced, math-specialized Llemma 34B \citep{llemma}, the latter having undergone additional training with Proof-Pile-2 \citep{llemma}. Of particular note, on the competition-grade MATH dataset, DeepSeekMath-Base exceeds the performance of all other open-source base models by more than 10\% in absolute terms, and even outperforms Minerva 540B \citep{minerva}—a proprietary model based on PaLM \citep{palm} that is approximately 77 times larger and specifically trained on mathematical corpora.

\paragraph{Mathematical Problem Solving with Tool Use} We assess program-aided mathematical problem solving for GSM8K and MATH by employing few-shot program-of-thought prompting \citep{pot,pal}. For each problem, models are tasked to devise a solution by writing a Python script that may leverage packages like \textit{math} and \textit{sympy} for complex calculations, and the result of script execution is used as the final answer. Results in Table \ref{tab:base_math_pot_proof} indicate that DeepSeekMath-Base 7B surpasses the previous leading model, Llemma 34B.

\paragraph{Formal Mathematics}

Automating formal proofs not only improves the rigor and reliability of mathematical reasoning but also enhances productivity, and has attracted increasing interest. We benchmark DeepSeekMath-Base 7B on the informal-to-formal proving task described by \citep{dsp_proof}, wherein the goal is to generate a formal proof from an informal statement, its formal rephrasing, and an informal argument. Evaluation is performed on miniF2F \citep{minif2f}, a suite for formalized Olympiad-style mathematical problems, prompting the model few-shot to generate Isabelle formal proofs for each problem. Following \cite{dsp_proof}, we use the models to construct proof outlines and then utilize the Sledgehammer automated theorem prover \citep{sledgehammer} to complete the remaining details. Table \ref{tab:base_math_pot_proof} shows that DeepSeekMath-Base 7B achieves robust performance in proof formalization.

\paragraph{Natural Language Understanding, Reasoning, and Code}

For assessments of language understanding, reasoning, and programming, we measure performance on MMLU \citep{mmlu} (natural language understanding), BBH \citep{bbh} (reasoning), HumanEval \citep{codex}, and MBPP \citep{mbpp} (code synthesis). As presented in Table \ref{tab:understanding_reasoning_code}, DeepSeekMath-Base 7B achieves considerable improvements over the earlier DeepSeek-Coder-Base-v1.5 \citep{deepseek-coder} in both MMLU and BBH, highlighting the beneficial effects of mathematical training on comprehension and reasoning tasks. Moreover, by incorporating code tokens in ongoing training, DeepSeekMath-Base 7B sustains the prior model’s results on the coding benchmarks. Overall, DeepSeekMath-Base 7B delivers a marked performance gain relative to the generalist Mistral 7B \citep{mistral} across all three evaluated reasoning and programming tasks.

\section{Supervised Fine-Tuning}

\subsection{SFT Data Curation}

We assembled a dataset for mathematical instruction tuning that encompasses problems in both English and Chinese, drawing from various mathematical disciplines and ranging in complexity. Each problem is matched with solutions expressed in chain-of-thought (CoT) \citep{cot}, program-of-thought (PoT) \citep{pot,pal}, as well as tool-integrated reasoning formats \citep{tora}. In total, the curated dataset consists of 776K training samples.

\begin{itemize}[topsep=0pt]
    \item \textbf{English mathematical datasets}: For English, we augmented GSM8K and MATH problems with solutions that incorporate external tool usage, and utilized a selected subset from MathInstruct \citep{MathInstruct} together with the training split of Lila-OOD \citep{lila}, where answers are generated using either CoT or PoT approaches. This collection spans a wide spectrum of mathematical areas, including algebra, probability, number theory, calculus, and geometry.
    \item \textbf{Chinese mathematical datasets}: For Chinese, we curated K-12 level mathematical problems across 76 different sub-domains, such as linear equations, each provided with solutions crafted in both CoT and tool-integrated styles.
\end{itemize}

\subsection{Training and Evaluating DeepSeekMath-Instruct 7B}

This section outlines the development of DeepSeekMath-Instruct 7B, which has been subject to instruction-based fine-tuning starting from the DeepSeekMath-Base model. During training, samples are randomly concatenated, ensuring the total token count does not exceed a 4K token context window. Model training proceeded for 500 steps, employing a batch size of 256, with a fixed learning rate set at 5e-5.

For evaluation, we assess the mathematical proficiency of the model, with and without tool assistance, across four quantitative reasoning benchmarks in both English and Chinese. Comparative performance is reported relative to top-performing contemporary models:

\begin{itemize}[topsep=0pt]
    \item \textbf{Closed-source models} considered are: (1) the GPT series, notably GPT-4 \citep{gpt4} and GPT-4 Code Interpreter~\footnote{\url{https://openai.com/blog/chatgpt-plugins\#\#code-interpreter}} as the strongest variants, (2) Gemini Ultra and Pro \citep{gemini}, (3) Inflection-2 \citep{inflection-2}, (4) Grok-1~\footnote{\url{https://x.ai/model-card}}, and select recent offerings from Chinese enterprises such as (5) Baichuan-3~\footnote{\url{https://www.baichuan-ai.com}} and (6) the newest GLM-4~\footnote{\url{https://open.bigmodel.cn/dev/api\#glm-4}} from the GLM model family \citep{glm}.
\end{itemize}

The majority of these models are designed for broad applications and have typically undergone several rounds of alignment.

\item \textbf{Open-source models} span: foundational models such as (1) DeepSeek-LLM-Chat 67B \citep{deepseek-llm}, (2) Qwen 72B \citep{qwen}, (3) SeaLLM-v2 7B \citep{seallm}, and (4) ChatGLM3 6B \citep{chatglm3}; as well as math-enhanced variants including (5) InternLM2-Math 20B~\footnote{\url{https://github.com/InternLM/InternLM-Math}}, which builds on InternLM2 through additional mathematics training and instruction fine-tuning; (6) Math-Shepherd-Mistral 7B, leveraging PPO training \citep{schulman2017proximal} on Mistral 7B \citep{mistral} paired with a process-supervised reward model; (7) WizardMath models \citep{wizardmath} that elevate mathematical reasoning in Mistral 7B and Llama-2 70B \citep{llama2} using evolve-instruct (an instruction tuning approach relying on instructions generated by AI) along with PPO training on datasets mainly comprising GSM8K and MATH; (8) MetaMath 70B \citep{metamath}, which is Llama-2 70B further refined with an expanded set of GSM8K and MATH problems; (9) ToRA 34B \cite{tora}, obtained by fine-tuning CodeLlama 34B for integrated mathematical problem-solving with external tools; (10) MAmmoTH 70B \citep{MathInstruct}, featuring Llama-2 70B adapted for mathematics via MathInstruct-based instruction tuning.
\end{itemize}

According to Table \ref{tab:sft_rl_math}, when evaluated without the aid of external tools, DeepSeekMath-Instruct 7B delivers robust step-wise reasoning performance. Remarkably, on the challenging MATH benchmark, this model outperforms all open-source alternatives and the vast majority of proprietary systems (including Inflection-2 and Gemini Pro) by a minimum margin of 9\% absolute, even when compared to models with significantly more parameters (such as Qwen 72B) or those refined with math-centric RL procedures (e.g., WizardMath-v1.1 7B). DeepSeekMath-Instruct’s results are on par with leading Chinese proprietary systems like GLM-4 and Baichuan-3 on MATH, though it does not yet reach the level of GPT-4 and Gemini Ultra.

Allowing models to combine natural language reasoning with programmatic tool use for problem resolution, DeepSeekMath-Instruct 7B achieves close to 60\% accuracy on MATH, outstripping all currently available open-source models. On additional benchmarks, its performance matches that of DeepSeek-LLM-Chat 67B—the preceding best-performing model, despite the latter's order-of-magnitude greater size.

\section{Reinforcement Learning}

\subsection{Group Relative Policy Optimization} The application of reinforcement learning (RL) following supervised fine-tuning (SFT) has proven beneficial in further advancing the mathematical reasoning capabilities of LLMs \citep{wang2023math,wizardmath}. In what follows, we detail Group Relative Policy Optimization (GRPO), our novel RL algorithm that is both efficient and effective.

\subsubsection{From PPO to GRPO} Proximal Policy Optimization (PPO) \citep{schulman2017proximal}, a widely-adopted actor-critic RL method, is extensively utilized during RL-based fine-tuning of LLMs \citep{ouyang2022training}. This approach updates model parameters by optimizing the surrogate objective function shown below:

\begin{equation}
\footnotesize
    \mathcal{J}_{PPO}(\theta) = \mathbb{E}{[q \sim P(Q), o \sim \pi_{\theta_{old}}(O|q)]} \frac{1}{|o|} \sum_{t=1}^{|o|} \min \left[ \frac{\pi_\theta(o_{t} | q, o_{<t})}{\pi_{\theta_{old}}(o_{t} | q, o_{<t})} A_{t}, \text{clip} \left( \frac{\pi_\theta(o_{t} | q, o_{<t})}{\pi_{\theta_{old}}(o_{t} | q, o_{<t})}, 1 - \varepsilon

Here, $\pi_{\theta}$ denotes the current policy model, while $\pi_{\theta_{old}}$ represents the preceding version of the policy. The variables $q$ and $o$ correspond to questions drawn from a dataset and output samples generated by the previous policy $\pi_{\theta_{old}}$, respectively. The hyper-parameter $\varepsilon$ relates to the clipping mechanism in PPO, which serves to ensure stable learning dynamics. The advantage term $A_t$ is computed using Generalized Advantage Estimation (GAE) \citep{gae}, which relies on rewards $\{r_{\ge t}\}$ and a value function $V_{\psi}$ learned alongside the policy. Consequently, PPO requires the simultaneous training of a value function in addition to the policy, and to avoid excessively exploiting the reward model, a standard practice is to append a per-token KL penalty—computed with respect to a reference model—to the reward at every token, following \citep{ouyang2022training}:

\begin{equation}
    r_{t} = r_\varphi(q, o_{\le t}) - \beta \log\frac{\pi_{\theta}(o_{t}|q, o_{<t})}{\pi_{ref}(o_{t}|q, o_{<t})},
\label{eq:PPO-reward}
\end{equation}

Here, $r_\varphi$ denotes the reward model, $\pi_{ref}$ is the reference model (typically initialized as the SFT model), and $\beta$ sets the strength of the KL regularization term.

Using a value function of similar scale as the policy model, as is customary in PPO, imposes notable memory and compute overheads. Further, the value function’s role in RL training is to act as a baseline within advantage calculations to reduce variance. In LLM settings, it is common to apply the reward model only to the final token, which introduces difficulties in training a token-level accurate value function. To overcome these challenges, we introduce Group Relative Policy Optimization (GRPO), as illustrated in Figure \ref{fig:grpo}. Unlike PPO, GRPO dispenses with separate value function approximation; instead, it employs the mean reward from a set of sampled outputs (each output responding to the same question) as the baseline. Concretely, for every question $q$, a group of responses $\{o_1, o_2, \cdots, o_G\}$ is generated via $\pi_{\theta_{old}}$. The policy is then optimized by maximizing the following objective:

\begin{equation}
\footnotesize
\begin{split}
    \mathcal{J}_{GRPO}(\theta) &= \mathbb{E}{[q \sim P(Q), \{o_i\}_{i=1}^G \sim \pi_{\theta_{old}}(O|q)]}  \\
    & \frac{1}{G}\sum_{i=1}^G\frac{1}{|o_i|} \sum_{t=1}^{|o_i|} \left\{ \min \left[ \frac{\pi_\theta(o_{i,t} | q, o_{i,<t})}{\pi_{\theta_{old}}(o_{i,t} | q, o_{i,<t})} \hat{A}_{i,t}, \text{clip} \left( \frac{\pi_\theta(o_{i,t} | q, o_{i,<t})}{\pi_{\theta_{old}}(o_{i,t} | q, o_{i,<t})}, 1 - \varepsilon, 1 + \varepsilon \right)  \hat{A}_{i,t} \right] - \beta \mathbb{D}_{KL}\left[\pi_{\theta} || \pi_{ref}\right]\right\} ,
\end{split}
\label{eq:GRPO-obj}
\end{equation}

Here, $\varepsilon$ and $\beta$ retain their roles as hyper-parameters, while $\hat{A}_{i,t}$ refers to the group-relative advantage, calculated solely based on the sampled outputs within each group—details of which follow in subsequent sections. The group-based computation of advantages in GRPO is well-suited to reward models, which are generally trained to make comparative judgments among responses to the same prompt. In addition, in contrast to incorporating a KL penalty within the reward as in PPO, GRPO applies KL regularization directly by adding the KL divergence term between the current policy and the reference policy into the objective, streamlining the evaluation of $\hat{A}_{i,t}$. Moreover, unlike the KL penalty applied in (\ref{eq:PPO-reward}), we estimate the KL divergence via the following unbiased estimator \

\begin{algorithm}[t]
  \small
  \caption{Iterative Group Relative Policy Optimization}
  \textbf{Input} initial policy model $\pi_{\theta_{\text{init}}}$; reward models $r_\varphi$; task prompts $\mathcal{D}$; 
  hyperparameters $\varepsilon$, $\beta$, $\mu$
  \begin{algorithmic}[1]
    \State Initialize policy model $\pi_\theta \leftarrow \pi_{\theta_{\text{init}}}$
    \For{iteration = 1, \dots, I}
       \State Assign current policy as reference model $\pi_{ref} \leftarrow \pi_{\theta}$
      \For{step = 1, \dots, M}
      \State Draw a mini-batch $\mathcal{D}_b$ from $\mathcal{D}$
      \State Set previous policy $\pi_{\theta_{old}} \leftarrow \pi_{\theta}$ 
      \State For each prompt $q \in \mathcal{D}_b$, sample $G$ outputs $\{o_i\}_{i=1}^G \sim \pi_{\theta_{old}} (\cdot \mid q) $
      \State Evaluate sampled outputs $o_i$ using the reward model $r_{\varphi}$ to obtain rewards $\{r_i\}_{i=1}^{G}$
      \State For every $t$-th token of $o_i$, estimate $\hat{A}_{i,t}$ using group relative advantage estimation
      \For{GRPO iteration = 1, \dots, $\mu$}
        \State Update policy $\pi_{\theta}$ by maximizing the GRPO objective (see Equation \ref{eq:GC-GRPO})
      \EndFor
    \EndFor
    \State Continuously update $r_\varphi$ through replay-based training.
    \EndFor 
  \end{algorithmic}
  \textbf{Output} $\pi_\theta$
  \label{alg:iter-grpo}
\end{algorithm}

\subsubsection{Outcome Supervision RL with GRPO} 
Formally, for each input $q$, a set of $G$ outputs $\{o_1, o_2, \cdots, o_G\}$ are sampled from the former policy $\pi_{\theta_{old}}$. Each output is assigned a reward by the reward model, producing a reward set $\mathbf{r} = \{r_1, r_2, \cdots, r_G\}$. To normalize the rewards, the group mean is subtracted and then divided by the group's standard deviation. In outcome supervision, the normalized reward associated with each output $o_i$ is provided at its end, and this normalized value is assigned as the advantage $\hat{A}_{i, t}$ for every token in that output, specifically, $\hat{A}_{i, t} = \widetilde{r}_i = \frac{r_i- {\rm mean}(\mathbf{r})}{{\rm std}(\mathbf{r})}$. Policy optimization proceeds by maximizing the objective introduced in equation (\ref{eq:GRPO-obj}).

\subsubsection{Process Supervision RL with GRPO} 
Outcome supervision limits feedback to the final step of each output, which may not provide adequate or granular guidance in tasks with complex mathematical reasoning. Building on \cite{wang2023math}, we implement process supervision, which delivers rewards at the completion of every intermediate reasoning step. For a question $q$ and $G$ sampled outputs $\{o_1, o_2, \cdots, o_G\}$, a process reward model assigns scores to each reasoning step, generating reward collections: $\mathbf{R} = \{ \{r_1^{index(1)},\cdots,r_1^{index(K_1)}\}, \cdots,  \{r_G^{index(1)},\cdots,r_G^{index(K_G)}\} \}$, where $index(j)$ identifies the end token of the $j$-th step and $K_i$ represents the step count for the $i$-th output. Each step reward is normalized by subtracting the average of all step rewards and dividing by their standard deviation, namely, $\widetilde{r}_i^{index(j)} = \frac{r_i^{index(j)} - {\rm mean(\mathbf{R})}}{{\rm std(\mathbf{R})}}$.
Under process supervision, the advantage assigned to the $t$-th token is computed as the sum of all normalized step rewards from the $t$-th token onward: $\hat{A}_{i, t} = \sum_{index

\subsubsection{Iterative RL with GRPO} During the progression of reinforcement learning training, the previous reward model might become inadequate to effectively guide the evolving policy model. To address this, we investigate the application of iterative RL with GRPO. As detailed in Algorithm \ref{alg:iter-grpo}, this method involves generating updated training datasets for the reward model using new samples from the current policy model. The previous reward model is then incrementally trained through a replay strategy that includes 10\% of past data. Subsequently, the policy model serves as the new reference, and its training continues using the updated reward model.

\subsection{Training and Evaluating DeepSeekMath-RL}

Our RL experiments utilize DeepSeekMath-Instruct 7B as the foundation. The RL training set comprises approximately 144K chain-of-thought-format questions sourced from GSM8K and MATH within the SFT data, purposefully omitting unrelated SFT questions to better isolate the effect of RL on benchmarks absent from the RL dataset. We construct reward model training data as outlined by \citep{wang2023math}. The initial reward model is developed from DeepSeekMath-Base 7B using a learning rate of 2e-5. For the GRPO stage, we use a learning rate of 1e-6 for the policy model and set the KL penalty coefficient to 0.04. For each problem, $64$ candidate solutions are sampled. We limit output sequences to a maximum of 1024 tokens and maintain a batch size of 1024 throughout training. The policy model undergoes only a single update after each round of exploration. For assessment, we evaluate DeepSeekMath-RL 7B on the same suite of benchmarks as DeepSeekMath-Instruct 7B. In this context, GSM8K and MATH, when solved via chain-of-thought reasoning, constitute in-domain tasks, while all other tested benchmarks are considered out-of-domain.

Table \ref{tab:sft_rl_math} presents results for both open- and closed-source models employing chain-of-thought and tool-integrated reasoning strategies across English and Chinese datasets. Our observations are as follows: 1) DeepSeekMath-RL 7B achieves chain-of-thought accuracies of 88.2\% on GSM8K and 51.7\% on MATH, outperforming every other open-source model within the 7B to 70B parameter range, as well as most closed-source counterparts. 2) Notably, DeepSeekMath-RL 7B is exclusively trained on chain-of-thought-style instruction tuning data from GSM8K and MATH, using DeepSeekMath-Instruct 7B as its base. Even with this limited training data, it consistently surpasses DeepSeekMath-Instruct 7B across evaluation metrics, highlighting the efficacy of reinforcement learning.

\section{Discussion}
In this section, we present insights gathered from both pre-training and reinforcement learning experiments.

\subsection{Lessons Learnt in Pre-Training}

We begin by detailing our experiences in pre-training. Unless stated otherwise, we follow the training setup described in Section \ref{sec:quality-policy}. Additionally, references to the DeepSeekMath Corpus in this section denote the 89B-token collection produced during the dataset's second version.

\subsubsection{Code Training Benefits Mathematical Reasoning}

A commonly held but insufficiently tested assumption is that exposure to code enhances reasoning capabilities. Our investigation contributes evidence supporting this claim, particularly for mathematical reasoning tasks, regardless of whether tool use is involved.

To examine the influence of code training on mathematical problem solving, we contrasted two paradigms: two-stage training and one-stage training.

\textbf{Two-Stage Training}

\begin{itemize}[topsep=0pt]
    \item \textbf{Code Training for 400B Tokens $\rightarrow$ Math Training for 150B Tokens}: We subject DeepSeek-LLM 1.3B to pretraining on 400B tokens of code data, subsequently followed by 150B tokens of mathematics data;
    \item \textbf{General Training for 400B Tokens $\rightarrow$ Math Training for 150B Tokens}: For baseline comparison, we substitute the code corpus in the initial phase with general domain data (randomly drawn from DeepSeek-AI’s expansive general text collection), aiming to elucidate the distinct impact of code data as opposed to general data on the model’s mathematical reasoning abilities.
\end{itemize}

\textbf{One-Stage Training}

\begin{itemize}[topsep=0pt]
    \item \textbf{Math Training for 150B Tokens}: DeepSeek-LLM 1.3B is directly exposed to 150B tokens of mathematics-specific training data;
    \item \textbf{Training on a mixture of 400B Code Tokens and 150B Math Tokens}: Math performance typically diminishes after code pretraining when followed by separate math training (due to catastrophic forgetting). To address this, we evaluate joint training, mixing code and math tokens within a single-stage process, to examine whether this approach sustains code capability and supports mathematical reasoning.
\end{itemize}

\paragraph{Results}
The quantitative outcomes of these experimental variations are presented in Table \ref{tab:code-math} and Table \ref{tab:code-math-general-eval}.

Pretraining on code yields marked advantages for program-augmented mathematical problem solving in both sequential and mixed-token regimens. According to Table \ref{tab:code-math}, two-stage training with an initial code-focused phase markedly strengthens DeepSeek-LLM’s effectiveness on GSM8K and MATH Python-based benchmarks. The addition of the subsequent math-focused stage further enhances this effect. Notably, in a one-stage setup, combining code and math tokens mitigates the catastrophic forgetting observed with a strict two-phase approach, resulting in improved performance both in code-related evaluations (Table \ref{tab:code-math-general-eval}) and program-aided math reasoning tasks (Table \ref{tab:code-math}).

Moreover, code-centric pretraining facilitates gains in mathematical reasoning tasks without recourse to tool usage. Within a two-stage structure, code-focused training independently produces moderate improvements, in addition to enhancing the downstream efficacy of subsequent math-targeted training—culminating in optimal results. On the contrary, simultaneous exposure to both code and math data in a one-stage procedure diminishes non-tool-based mathematical reasoning, potentially because DeepSeek-LLM 1.3B's restricted parameter count limits its ability to absorb both modalities concurrently.

\subsubsection{ArXiv Papers Appear Unhelpful for Advancing Mathematical Reasoning}

ArXiv papers are frequently utilized as a fundamental source for mathematical pre-training datasets \citep{minerva,gpt-f,llemma,mathpile}. Despite their prevalence, a thorough evaluation of their specific contribution to mathematical reasoning remains largely unexplored. Somewhat unexpectedly, our experimental results indicate that including arXiv papers yields little to no improvement in mathematical reasoning ability. Our study examines a range of model scales, specifically DeepSeek-LLM 1.3B and DeepSeek-Coder-Base-v1.5 7B \citep{deepseek-coder}, leveraging arXiv datasets that have undergone different processing strategies:

\begin{itemize}[topsep=0pt]
    \item \textbf{MathPile} \citep{mathpile}: an 8.9B-token dataset assembled through cleaning and heuristic filtering procedures, consisting of over 85\% scientific papers from arXiv;
    \item \textbf{ArXiv-RedPajama} \citep{redpajama}: the full arXiv LaTeX collection after removing preambles, comments, macros, and bibliographies, amassing 28.0B tokens.
\end{itemize}

Our methodology involved separately pre-training DeepSeek-LLM 1.3B for 150B tokens and DeepSeek-Coder-Base-v1.5 7B for 40B tokens on each version of the arXiv data. The outcomes demonstrate that training exclusively on arXiv-based corpora does not facilitate, and may in fact hinder, mathematical reasoning. This observation holds consistently across various mathematical benchmarks spanning different levels of complexity featured in this work, such as GSM8K and MATH for quantitative reasoning (see Table \ref{tab:arxiv-cot}), MMLU-STEM for multiple-choice testing (see Table \ref{tab:arxiv-cot}), and miniF2F for formal mathematics (refer to Table \ref{tab:arxiv_atp}).

Nevertheless, several caveats remain, warranting cautious interpretation of these findings. In particular, we did not explore:

\begin{itemize}[topsep=0pt]
    \item The potential influence of arXiv-derived data on specialized mathematical subtasks absent from our benchmarks, such as translating formal proofs or statements into informal text;
    \item The effects when arXiv tokens are incorporated alongside alternative data types;
    \item The possibility that benefits from arXiv sources might become apparent at increased model scales.
\end{itemize}

Accordingly, these questions necessitate additional research, which we intend to pursue in the future.

\subsection{Reflections on Reinforcement Learning}
\subsubsection{Toward a Unified Framework} In this portion, we introduce a consolidated perspective to examine multiple training techniques, including SFT, RFT, DPO, PPO, and GRPO. We supplement this framework with experiments analyzing factors within this unified approach. Generally, for a given training technique, the gradient with respect to parameters $\theta$ may be represented as:

\begin{equation}
    \nabla_{\theta}\mathcal{J}_{\textcolor{red}{\mathcal{A}}}(\theta) = \mathbb{E}[\underbrace{(q,o) \sim \textcolor{red}{\mathcal{D}}}_{Data \ Source}]\left( \frac{1}{|o|} \sum_{t=1}^{|o|}  \underbrace{GC_{{\mathcal{A}}}(q, o, t, \textcolor{red}{\pi_{{rf}}})}_{Gradient \ Coefficient}  \nabla_{\theta}\log \pi_{\theta}(o_t | q, o_{<t})\right).
\label{eq:objective}
\end{equation}

This formulation comprises three principal elements: 1) \textit{Data Source $\mathcal{D}$}, which specifies the set of examples used during training; 2) \textit{Reward Function $\pi_{{rf}}$}, responsible for providing the learning signal; and 3) \textit{Algorithm $\mathcal{A}$}, which integrates both the training examples and reward feedback into the gradient coefficient $GC$, thereby controlling the adjustment strength—either punitive or reinforcing—applied to each data instance. Under this generalized scheme, we review several archetypal approaches:

\begin{itemize}
    \item  \textbf{Supervised Fine-tuning (SFT)}: SFT entails adapting a pretrained model by further training on a curated human-annotated SFT dataset.
    \item \textbf{Rejection Sampling Fine-tuning (RFT)}: RFT advances the SFT model through additional training on a filtered subset of responses generated by the SFT model for SFT queries; this filtering mechanism retains only responses that meet correctness criteria.
    \item \textbf{Direct Preference Optimization (DPO)}: DPO further polishes the SFT model by fine-tuning on an expanded set of candidate completions derived from the SFT model, optimizing according to the pairwise DPO objective.
    \item \textbf{Online Rejection Sampling Fine-tuning (Online RFT)}: Distinct from conventional RFT, Online RFT starts with the SFT model as the initial policy and iteratively fine-tunes on dynamically sampled augmented outputs from the up-to-date policy during training.
    \item \textbf{PPO/GRPO}: PPO/GRPO begin with the SFT model to initialize the policy and reinforce it by sampling outputs from the continuously updated policy, applying reinforcement learning.
\end{itemize}

The key elements of these approaches are outlined in Table \ref{tab:data_policy}. For an expanded explanation of the derivations, see Appendix \ref{app:analysis-rl}.

\paragraph{Insights on Data Source} We categorize data sources into two primary types: online sampling and offline sampling. Online sampling refers to data generated from the current training policy model’s exploration, whereas offline sampling involves using data sampled from the initial SFT model. Methods such as RFT and DPO adopt offline sampling, whereas Online RFT and GRPO utilize the online sampling approach.

As depicted in Figure \ref{fig:rl-analysis}, our results indicate that Online RFT yields notably better performance than RFT on both evaluation benchmarks. In the early phases of training, the performance of Online RFT closely matches that of RFT; however, as training proceeds, Online RFT gains a clear lead, emphasizing the effectiveness of online data collection. This pattern is reasonable: early in training, the actor closely mirrors the SFT model, leading to limited variability in sampled data. As training progresses, samples from the actor model increasingly diverge from the SFT model, at which point real-time online data sampling brings substantial benefits.

\paragraph{Insights on Gradient Coefficient} The reward signal, after being processed into a gradient coefficient, guides updates to the model parameters. We distinguish between two types of reward functions in our experiments: `Rule' and `Model.' The Rule approach evaluates responses based on their factual correctness, whereas the Model method utilizes a reward model trained on data labeled by the Rule mechanism to score responses. Equations \ref{eq:GC-RFT} and \ref{eq:GC-GRPO} illustrate a central distinction: GRPO dynamically modifies its gradient coefficient according to the reward model’s scores, enabling differential rewarding or penalization based on response quality. Online RFT, in contrast, does not have this capacity—it applies consistent reinforcement to all correct responses and does not impose penalties on incorrect ones.

As illustrated in Figure \ref{fig:rl-analysis}, GRPO outperforms online RFT, underscoring the effectiveness of adjusting positive and negative gradient coefficients. Moreover, GRPO+PS yields better outcomes than GRPO+OS, demonstrating the value of leveraging step-aware, fine-grained gradient coefficients. We further investigate the effects of iterative RL by implementing two iterations in our experiments. Figure \ref{fig:iter-rl} reveals that iterative RL substantially boosts performance, with the first iteration contributing the most significant improvement.

\subsubsection{Why RL Works?} In this study, reinforcement learning is applied using a subset of instruction tuning data, leading to notable improvements over the baseline instruction tuning model. To elucidate the mechanism behind RL’s success, we assess the Pass@K and Maj@K metrics of both Instruct and RL models across two benchmarks. As depicted in Figure \ref{fig:maj-pass}, RL notably elevates the Maj@K metric, whereas Pass@K remains largely unchanged. These results suggest that RL strengthens the overall robustness of the output distribution, namely, \textbf{the observed gains arise mainly from amplifying the likelihood of correct answers within the TopK outputs, rather than from intrinsic enhancements in model ability.} Correspondingly, \citep{wang2023making} identified a \textbf{misalignment problem} in reasoning tasks within SFT models and showed that performance in these tasks can be advanced by employing various preference alignment strategies \citep{yuan2023rrhf,song2023preference,wang2023making}.

\subsubsection{How to Achieve More Effective RL?}
\label{sec:effec_rl}
Our results indicate that RL can achieve strong performance on mathematical reasoning problems. Additionally, we introduce a unified conceptual framework that encapsulates multiple representative training approaches under the lens of direct or simplified RL methodologies. As encapsulated by Equation \ref{eq:objective}, the essential components consist of Data Source, Algorithm, and Reward Function. We discuss prospective directions concerning each of these components.

\paragraph{Data Source} The data source serves as the foundational element for all training procedures. Within RL, it specifically pertains to unlabeled questions paired with sampled outputs from the policy model. In this work, we restrict the data source to questions used during instruction tuning, relying on simple nucleus sampling to generate responses. We speculate that this limitation is a likely reason why RL in our setting primarily augments the Maj@K metric. In future research, we intend to extend our RL pipeline to incorporate out-of-distribution question prompts and adopt \textbf{more sophisticated sampling (decoding) strategies}, including those leveraging tree-search techniques \citep{yao2023tree}. Furthermore, advancements in \textbf{efficient inference techniques} \citep{xia-etal-2023-speculative,leviathan2023fast,kwon2023efficient,xia2024unlocking}, which dictate the exploration efficiency of policy models, are anticipated to be critically important.

\paragraph{Algorithms} Algorithms utilize both the input data and the received reward signals to compute the gradients necessary for updating the model parameters. In accordance with Equation \ref{eq:objective}, prevailing approaches tend to place significant \textbf{TRUST} in the reward function, using it to either promote or penalize the conditional probability of specific tokens. Nonetheless, the reliability of the reward signal cannot be universally guaranteed, particularly when faced with highly intricate tasks. As an illustration, even the PRM800K datasets \citep{lightman2023let}, curated with considerable effort by skilled annotators, still exhibit an estimated error rate of around 20\% in annotations\footnote{\url{https://github.com/openai/prm800k/issues/12\#issuecomment-1728491852}}. Therefore, there is a need to investigate reinforcement learning algorithms that exhibit resilience to noisy or imperfect reward feedback. We argue that the advancement and adoption of \textbf{WEAK-TO-STRONG} \citep{burns2023weak} alignment paradigms have the potential to revolutionize the design of learning algorithms at a foundational level.

\paragraph{Reward Function} The reward function acts as the fundamental provider of feedback for training. Within the reinforcement learning paradigm, this function is typically operationalized through a neural reward model. We identify three pivotal avenues for the evolution of reward models: 1) \textbf{Enhancing the reward model's generalization capacity.} The reward model should be capable of robustly generalizing to unseen inputs and sophisticated outputs from advanced decoding, lest reinforcement learning merely conserves the status quo of LLM distributions without driving substantive improvement; 2) \textbf{Quantifying the uncertainty inherent to reward models.} Explicitly modeling this uncertainty may facilitate more seamless integration between imperfect reward models and weak-to-strong learning mechanisms; 3) \textbf{Developing efficient, high-fidelity process reward models} that deliver nuanced and detailed training signals throughout reasoning chains \citep{lightman2023let,wang2023math}.

\section{Conclusion, Limitation, and Future Work}

We introduce DeepSeekMath, which sets a new benchmark for open-source models on competition-level MATH tasks, narrowing the performance gap with leading closed-source systems. The model is initialized using DeepSeek-Coder-v1.5 7B and is continually trained on a corpus of 500B tokens, prominently featuring 120B mathematical tokens extracted from Common Crawl. Our comprehensive ablation studies reveal that mathematical content from web resources provides a rich and valuable dataset, whereas arXiv data does not yield as much benefit as anticipated. In this work, we propose Group Relative Policy Optimization (GRPO), a modification of Proximal Policy Optimization (PPO), which substantially enhances mathematical reasoning ability while using less memory. Experimental results indicate that GRPO remains beneficial even when DeepSeekMath-Instruct 7B achieves strong results on various benchmarks. Additionally, we propose a unified framework to contextualize different approaches and highlight several avenues for advancing reinforcement learning effectiveness.

While DeepSeekMath attains outstanding results on quantitative reasoning benchmarks, its proficiency in geometry and theorem proving lags behind closed models. For example, during preliminary testing, DeepSeekMath struggled with problems involving triangles and ellipses, potentially revealing selection biases in both pre-training and fine-tuning datasets. Furthermore, due to the model’s size, DeepSeekMath demonstrates inferior few-shot performance compared to GPT-4; the latter displays marked gains with few-shot inputs, whereas DeepSeekMath performs similarly in zero-shot and few-shot scenarios. Looking ahead, we aim to refine our data selection process to build higher-quality pre-training datasets. We also plan to pursue the research directions identified in Section \ref{sec:effec_rl} to foster more efficient reinforcement learning for LLMs.

\end{document}